\documentclass[french]{book}
\usepackage{teach}
\usepackage{verbatim}
\usepackage{amsmath,amsfonts,amssymb}
\usepackage{pgfplots}
\usepackage{mwe}
\usepackage{stmaryrd}
\usepackage{yhmath} 
\usepackage{pstricks,pst-plot,pst-text,pst-tree,pst-eucl}
\newcommand{\bsyc}{\left\{\begin{array}{rcl}} % syst�e serr�sur le �al
\newcommand{\esyc}{\end{array}\right.}


\newcommand{\ofr}[2]{\dfrac{#1}{#2}}
\newcommand{\arc}[1]{\wideparen{#1}}
\RequirePackage{graphicx}
\RequirePackage{ifpdf}
 \ifpdf
   \DeclareGraphicsRule{*}{mps}{*}{}
 \fi

\newcommand{\sysd}[2]{%
       \left\{
       \begin{aligned}
          #1\\
          #2\\
       \end{aligned}
       \right.%
       }
  \usepackage{fancyhdr}
  \usepackage{tkz-tab}
  \usepackage{amsmath}
  \usepackage{amssymb}
  \usepackage{amsfonts}
  \usepackage{amsthm}
  \usepackage{mathrsfs}

  \usepackage{array}
  \usepackage{multicol}
  \setlength{\multicolsep}{3pt}

  \usepackage{graphicx}
  \graphicspath{{Figures/}}
  \usepackage{pstricks,pst-plot,pst-text,pst-tree,pst-eucl}
  \usepackage[all]{xy}

  \usepackage{fancybox}
  \usepackage{color}

%  \usepackage[frenchb]{babel}
  \usepackage{hyperref}
\usepackage{textcomp}
\usepackage{mathtools}
\usepackage{subfig}
\pgfplotsset{compat=newest}
\newcommand{\I}{i}
\usepackage[all]{xy}
%\usepackage{geometry}
%\geometry{top=1cm, bottom=1cm, left=2cm, right=2cm}
\usepackage[utf8]{inputenc}
\usepackage[french]{babel}
%\usepackage[bottom]{footmisc}
\redefinecolor{thm}{title}{bgcolor}{alizarin}
\redefinecolor{prop}{title}{bgcolor}{meatbrown}
\redefinecolor{defn}{title}{bgcolor}{olivine}
\definecolor{alizarin}{rgb}{0.82, 0.1, 0.26}
\definecolor{meatbrown}{rgb}{0.9, 0.72, 0.23}
\definecolor{olivine}{rgb}{0.6, 0.73, 0.45}
\usepackage{mathrsfs}
%\usepackage{preambule}
%%
\newcommand{\R}{\mathbb {R}}
\newcommand{\bbr}{\mathbb {R}}
\newcommand{\bbc}{\mathbb {C}}
\newcommand{\bbz}{\mathbb {Z}}
%\newcommand{\C}{\mathds {C}}
\newcommand{\N}{\mathbb {N}}
\newcommand{\Z}{\mathbb {Z}}
\newcommand{\Q}{\mathbb {Q}}
\newcommand{\D}{\mathbb {D}}
\newcommand{\HRule}{\rule{\linewidth}{0.5mm}}
%%
\newcommand{\se}{\geq}
\newcommand{\ie}{\leq}
%\newcommand{\imath}{i}
\newcommand{\tm}{\times}
\newcommand{\ab}[1]{\vert #1 \vert}
\newcommand{\nov}[1]{\Vert \overrightarrow{#1} \Vert}
%\newcommand{\ell}{l}
%\newcommand{\ell'}{l'}
%%
\newcommand{\barre}[1]{\overline{#1}}
\newcommand{\ds}{\displaystyle}
\newcommand{\limn}{\lim_{n\rightarrow +\infty}}
\newcommand{\limo}[1][x]{\ds\lim_{#1\rightarrow 0}}
\newcommand{\limite}[2][x]{\ds\lim_{#1\rightarrow #2}}
\newcommand{\limpinf}[1][x]{\ds\lim_{#1\rightarrow +\infty}}
\newcommand{\limminf}[1][x]{\ds\lim_{#1\rightarrow -\infty}}
\newcommand{\limiteg}[2][x]{\ds\lim_{#1\xrightarrow{<}#2}}
\newcommand{\limited}[2][x]{\ds\lim_{#1\xrightarrow{>}#2}}
%%
\newcommand{\vect}[1]{\overrightarrow{#1}}
\newcommand{\oij}{(\textit{O},{\overrightarrow{i}},{\overrightarrow{j}})}
%
\newcommand{\co}[2]{C_#2^#1}
\newcommand{\ud}[1]{\dfrac{#1}{2}}
\newcommand{\cp}[2]{%
        \begin{pmatrix}
        #1\\
        #2
        \end{pmatrix}%
        }
%
\newcommand{\calb}{\mathscr{B}}
\newcommand{\calc}{\mathscr{C}}
\newcommand{\cald}{\mathscr{D}}
\newcommand{\cale}{\mathscr{E}}
\newcommand{\calf}{\mathscr{F}}
\newcommand{\calp}{\mathscr{P}}
\newcommand{\cals}{\mathscr{S}}
\newcommand{\calt}{\mathscr{T}}
\newcommand{\calr}{\mathscr{R}}
\newcommand{\cali}{\mathscr{I}}
\newcommand{\RR}{\calr }
\newcommand{\CR}{\calc}
\newcommand{\IR}{\cali }
\newcommand{\fr}[2]{\dfrac{#1}{#2}}
\newcommand{\ve}[1]{\overrightarrow{#1}}
\newcommand{\pa}[1]{(#1)}
\newcommand{\ps}[2]{\overrightarrow{#1}.\overrightarrow{#2}}
%\newcommand{\bbr}{\mathbb{R}}
\newcommand{\al}{\alpha}
\newcommand{\la}{\lambda}
\newcommand{\DR}{\cald}
\newcommand{\PR}{\calp}
\newcommand{\cacher}[1]{#1}
\newcommand{\cacheg}[1]{#1}
\newcommand{\cacheb}[1]{#1}
\newcommand{\nor}[1]{ \Vert #1 \Vert}
\newcommand{\abv}[1]{\vert #1 \vert}

\newenvironment{arrayl}[1][rcl]%
	{\setlength{\arraycolsep}{1.5pt}
	 $\begin{array}[t]{#1}}%
	{\end{array}$}%
%\newcommand{\coordpp}[3]{\begin{pmatrix}
%
\newcommand{\suite}[1][u]{\left( #1_{n} \right)_{n\in\N}}
\newcommand{\suitep}[1][u]{\left( #1_{n} \right)_{n\in\N^{*}}}
\usetikzlibrary{arrows}
\newcommand{\ssi}{\Longleftrightarrow}
\newcommand{\poly}[1][a]{#1_0+#1_1x+\cdots+#1_nx^n}
\newcommand{\coordpp}[3]{\begin{pmatrix}
#1\\
#2\\
#3
\end{pmatrix}}
\newcommand{\anglevec}[2]{(\overrightarrow{#1};\overrightarrow{#2})}
\newcommand{\un}{\suite}
\newcommand{\eit}{e^{it}}
\newcommand{\eimt}{e^{imt}}
\newcommand{\vn}{\suite[v]}
%\newcommand{\pa}[1]{(#1)}
\newcommand{\cro}[1]{[#1]}
\newcommand{\abs}[1]{\vert #1 \vert}

\newcommand{\Syst}[2]{\left\{\begin{array}{lcccc} #1 \\ #2 \end{array}\right.}
\renewcommand{\arraystretch}{1.2} 
\newcolumntype{C}[1]{>{\centering\arraybackslash}p{#1cm}}
\newcommand{\Rep}{(O;\ve{i};\ve{j})}
\newcommand{\C}[2]{C_{#1}^{#2}}
\newcommand{\dx}{dx}

\begin{document}
\tableofcontents
\chapter{Arithmétique}

\section{Divisibilité dans $\mathbb{Z}$ et congruences}

 \begin{defn}
Soit $a$ et $b$ deux entiers relatifs.
On dit que $a$ \underline{divise} $b$ s'il existe un entier $k$ tel que $b=a\times k$.
Lorsque $a$ divise $b$, on note $a\mid b$ 

On dit aussi que: 
\begin{tabular}[t]{lll}$a$ est un \textbf{diviseur} de $b$\\$b$ est \textbf{divisible} par $a$\\$b$ est un    \textbf{multiple} de $a$
\end{tabular}

\end{defn}

 \begin{rem}

$0$ est un multiple de tout entier: pour tout $n\in \mathbb{Z}, 0 \times n=0$


 \end{rem}

\begin{prop}
Soit $a$ et $b$ deux entiers

a) Si $a \mid b$ et $b \neq 0 $ alors $ \,\vert a \vert \leqslant \vert b \vert$

b) Tout entier $b$ \textbf{non nul }a un nombre fini de diviseurs.

\end{prop}

 \begin{demo}
a) Comme $a \mid b$, on peut écrire $b=k \times a$ avec $k\in \mathbb{Z}$ donc:

$\vert b \vert=\vert k \times a \vert = \vert k\vert \times \vert a\vert$
mais $b\neq 0$ donc $k \neq 0$ et $\vert k\vert \geq1$ et par conséquent:$\vert a\vert\leq\vert b\vert$

b) Soit $b$ un entier non nul. Si $a$ est un diviseur de $b$ on a vu que: $\vert a\vert\leq\vert b\vert$
donc $-\vert b \vert\leq a\leq \vert b \vert$ et $a$ peut prendre au maximum $2 \times \vert b \vert$ valeurs.

En revanche, 0 a une infinité de diviseurs car tous les entiers divisent 0.

\end{demo}


 \begin{prop} 
Soit $a, b$ et $c$ trois entiers.

a) Si $a\mid b$ et $b\mid c$ alors $a\mid c$

b) $c\mid a$ si et seulement si $c\mid -a$

{\small (L'ensemble des diviseurs de $a$ est égal à l'ensemble des diviseurs de $-a$)}



c)Si $c\mid a$ alors $c\mid ab$

d) Si $c\mid a$ et $c\mid b$ alors: \begin{tabular}[t]{lll}1) $c\mid a+b$\\2) $c\mid a-b$\\3) $c\mid au+bv$ avec $u$ et $v$ entiers
\end{tabular}

e) Soit $a$ et $b$ non nuls. Si $b\mid a$ et $a \mid b$ alors $a=b$ ou $a= -b$

\end{prop}

 \begin{demo}

b) Si $c\mid a$, il existe un entier $k$ tel que: $a=k \times c$ et donc $-a=\left( -k\right) \times c$ soit $c\mid-a$

Si $c\mid -a$, il existe un entier $k'$ tel que: $-a=k' \times c$ et donc $a=\left( -k'\right) \times c$ soit $c\mid a$

\end{demo}






\begin{thm}
Soit $a$ un entier et $b$ un naturel non nul.

Il existe un unique entier $q$ et un unique entier $r$ tels que:
\begin{center}
$a=b \times q+r$ avec $0\leq r < b$
\end{center}

\end{thm}

\begin{demo}
On considère les multiples de $b: ...-2b, -b, 0, b, 2b, ...$

L'entier $a$ est:

* soit multiple de $b$ et dans ce cas il existe un unique entier $q$ tel que: $a=b \times q$

* soit compris strictement entre deux multiples de $b$ et donc il existe un unique entier $q$ tel que:

\hspace{3cm}$bq < a < b\left( q+1\right) $

Si on note $r=a-bq$ on a bien, dans les deux cas, $a=b \times q+r$ avec $0\leq r < b$

\end{demo}


\begin{rem}

$a$ est le dividende,
$b$ est le diviseur,
$q$ est le quotient et
$r$ est le reste
\end{rem}

\begin{exemple}
Déterminer le quotient et le reste de la division euclidienne de $a$ par $b$ avec:

a) $a=325, b=7$

b) $a=-113, b=7$


Réponses: \\

a) $325=7 \times 46+3$

b) $113=7 \times 16 +1$ d'où $-113=-7 \times 16 -1$ puis $-113=-7 \times 17+6$, $q=-17, r=6$


\end{exemple}






\begin{defn}
Soit $a$ et $b$ deux entiers et $n$ un entier naturel non nul.

On dit que $a$ est \textbf{congru à} $b$ \textbf{modulo} $n$ si $a$ et $b$ ont le même reste dans la division euclidienne par $n$.

\end{defn}

\begin{rem}

$a$ est congru à $b$ modulo $n$ se note au choix:
$a \equiv b \left( n\right) $ ou $a \equiv b \left[ n\right] $ ou $a \equiv b$ mod $ n$
 
 \end{rem}
 
\begin{exemple}
86 et 23 ont pour reste 2 dans la division euclidienne par 7 donc $86\equiv 23 \left( 7\right) $

\end{exemple}

\begin{thm}
Soit $a$, $b$, $c$ et $r$ entiers et $n$ un entier naturel non nul.

1) $a \equiv b \left( n\right) $ si et seulement si $a-b$ est un multiple de $n$

2) Si $a \equiv r \left( n\right) $ et si $0 \leq r < n$ alors $r$ est le reste de la division de $a$ par $n$.

3) Si $a \equiv b \left( n\right) $ et  $b \equiv c \left( n\right) $ alors $a \equiv c \left( n\right) $

\end{thm}



\begin{demo}[ $a \equiv b \left( n\right) $ si et seulement si $a-b$ est un multiple de $n$]

  $\bigstar$ Si $a \equiv b \left( n\right) $ alors la division  par $n$ donne: $a=nq+r$ et $b=nq'+r$ donc $a-b=n(q-q')$: $n$ divise $a-b$
  
 $\bigstar$ Si $n$ divise $a-b$ il existe $k \in\mathbb{Z}$ tel que $a-b=nk$ soit $a=b+nk$
  
  La division de $a$ par $n$ se traduit par: $a=nq+r$ avec $0\leq r <n$. On a alors $b+nk=nq+r$ donc $b=n(q-k)+r$ et $0\leq r < n$ qui se traduit par: $a$ et $b$ ont le même reste dans la division par $n$

\end{demo}


\begin{prop}[Congruences et opérations]
Soit $a$, $b$, $c$ et $d$ entiers et $n$ un entier naturel non nul.

1) Si $a \equiv b \left( n\right) $ alors $ac \equiv bc \left( n\right) $

2) Si $a \equiv b \left( n\right) $ et $c \equiv d \left( n\right) $ alors:

\begin{tabular}[t]{lll}1) $a+c \equiv b+d \left( n\right) $\\2) $a-c \equiv b-d \left( n\right) $\\3) $ac \equiv bd \left( n\right) $
\end{tabular}

3) Si $a \equiv b \left( n\right) $ alors, pour tout naturel $p$, on a: $a^{p} \equiv b^{p} \left( n\right) $

\end{prop}

\section{PGCD}




\begin{defn}

Soit $a$ et $b$ deux entiers. On note $\Delta\left( a;b\right) $ l'ensemble des diviseurs communs à $a$ et $b$.
Si $b\mid a$ alors $\Delta\left( a;b\right) $ est égal à l'ensemble des diviseurs de $b$, qu'on note $\Delta\left( b\right) $

\end{defn}

\begin{prop}
Si $a, b, q$ et $r$ sont des entiers tels que: $a=bq+r$ alors $\Delta\left( a;b\right)=\Delta\left( b;r\right)  $
 \end{prop}
 
\begin{demo}

Si $d\in \Delta\left( a;b\right)$ alors $d\mid r$ car $r=a-bq$ donc $d\in \Delta\left( b;r\right)$

Si $d\in \Delta\left( b;r\right)$ alors $d\mid a$ car $a=bq+r$ donc $d\in \Delta\left( a;b\right)$


Remarque: cas particulier\\

Si $a\in \mathbb{Z}, b\in\mathbb{N}^{*}$ et $a=bq+r, 0 \leq r < b $ est la division euclidienne de $a$ par $b$ alors: $\Delta\left( a;b\right)=\Delta\left( b;r\right)  $


\end{demo}

\begin{defn}

Si $a$ et $b$ sont deux entiers non tous les deux nuls alors $ \Delta\left( a;b\right)$ est un ensemble fini et on appelle PGCD de $a$ et $b$ le \textbf{plus grand diviseur commun} de $a$ et $b$.
\end{defn}

\begin{rem}
pgcd$(a;b)$ ou $a\wedge b$
\end{rem}

\begin{prop}
Soit $a$, $b$ et $r$ entiers non nuls.

1) pgcd$(a;b)\geq 1$

2) pgcd$(a;b)$ = pgcd$(b;a)$

3) pgcd$(a;a)=\vert a\vert$

4) pgcd$(a;b)$ =pgcd$(-a;b)$ =pgcd$(a;-b)$

5) Si $b\mid a$, pgcd$(a;b)=\vert b \vert$

6) Si $a=bq+r$  alors,  pgcd$(a;b)$ =pgcd$(b;r)$ 


\end{prop}

\begin{exemple}[Principe de l'algorithme d'Euclide pour calculer le pgcd]

Soit $a$ et $b$ deux entiers naturels non nuls tels que $b$ ne divise pas $a$. Alors pgcd$(a;b)$ est le \textbf{dernier reste non nul} de la suite des divisions de l'algorithme d'Euclide.

Si $r_{1},r_{2},...,r_{n}$ est la suite de restes non nuls, on a:

$ \Delta\left( a;b\right)= \Delta\left( b;r_{1}\right)=... \Delta\left( r_{n-1};r_{n}\right)= \Delta\left( r_{n}\right)$


\end{exemple}

\begin{exemple}

Calculer pgcd$(8820;3150)$

\underline{Solution}
$8820=3150\times 2+2520$

$3150=2520\times 1+630$

$2520=630\times 4+0$ 
donc pgcd$(8820;3150)=630$

\end{exemple}

\begin{prop}

Soit $a$ et $b$ deux entiers non nuls et $d= pgcd(a;b)$

1) L'ensemble des diviseurs communs à $a$ et $b$ est l'ensemble des diviseurs de $d$

2) Il existe des entiers $u$ et $v$ tels que: $d=au+bv$

3) Pour  $k \in\mathbb{N}^{*}, pgcd(ka;kb)=k\times pgcd(a;b)$

\end{prop}

\begin{demo}

1)$\Delta(a;b)=\Delta(d)$ d'après l'algorithme d'Euclide.

2) Soit $E=\left\lbrace au+bv \text{ avec } u,v\in \mathbb{Z}\right\rbrace $

Si $u=1,v=0$, on voit que $a\in E$
Si $u=-1,v=0$, on voit que $-a\in E$
Si $u=0,v=1$, on voit que $b\in E$

$E$ contient au moins un entier strictement positif. Soit $\delta$ le plus petit d'entre eux;
il existe $u_{0}$ et $v_{0}$ entiers tels que: $\delta=au_{0}+bv_{0}$

Soit $au+bv$ un élément quelconque de $E$ La division euclidienne de cet élément par $\delta$ donne
  $au+bv=\delta q+r, 0 \leq r < \delta $ d'où: $r=au+bv-q \left(au_{0}+bv_{0}\right)=a \left( u-qu_{0}\right) +b\left( v-qv_{0}\right) $
 
 
 
 Par suite, $r\in E$ et $0 \leq r < \delta$, mais d'après la définition de $\delta$: $r=0$ donc \textbf{tout élément de E est multiple de } $\delta$.
 
 $\delta\mid a$ et $\delta\mid b$ donc $\delta\mid d=pgcd(a;b)$.
 
 Réciproquement, $d\mid a$ et $d\mid b$ donc aussi $d\mid au_{0}+bv_{0}$ c'est-à-dire, $d\mid\delta$. Conclusion: $d=\delta$
 
 
\end{demo}




\section{Entiers premiers entre eux}


\begin{defn}

Soit $a$ et $b$ deux entiers non nuls. On dit que $a$ et $b$ sont \textbf{premiers entre eux} si pgcd$(a,b)=1$

\end{defn} 
 
\begin{thm}

Soit $a,b$ et $d$ des entiers non nuls. 

pgcd$(a,b)=d$ si et seulement si, il existe deux entiers non nuls $a'$ et $b'$ tels que: $a=da',b=db'$ et $a'\wedge b'=1$
 
\end{thm}

\begin{demo}

  Si $d=$ pgcd$(a,b), d\mid a, d\mid b$ donc il existe $a', b' \in \mathbb{Z}^{*}$ tels que: $a=da', b=db'$
  
  Soit $r=$ pgcd$(a',b')$ alors pgcd$(da',db')=dr=d$ donc $dr=d$ avec $d\neq 0$ donc $r=1$ et $a'\wedge b'=1$
  
  Réciproquement, si $a=da', b=db'$ et $a'\wedge b'=1$ alors pgcd$(da',db')=d \times$pgcd$(a',b')=d$
 
\end{demo}


\begin{thm}[Théorème de Bézout]

Deux entiers non nuls $a$ et $b$ sont premiers entre eux si et seulement si, il existe des entiers $u$ et $v$ tels que: $au+bv=1$
\end{thm}

\begin{demo}

Si $a$ et $b$ premiers entre eux, pgcd$(a,b)=1$ et donc il existe deux entiers $u$ et $v$ tels que $au+bv=1$

Réciproquement, si $au+bv=1$ et $d\mid a$ et $d\mid b$ alors $d\mid au+bv$ donc $d$ divise 1. Les seuls diviseurs communs à $a$ et $b$ sont $-1$ et $1$ et $a \wedge b=1$

\end{demo}

\begin{cor}
Soit a,b deux entiers tel que $pgcd(a,b)=d$ alors il existe deux entiers $u$ et $v$ tels que,

 $$ua+vb=d$$
\end{cor}

\begin{demo}
Soit a,b deux entiers tel que $pgcd(a,b)=d$ alors $pgcd(\frac{a}{d},\frac{b}{d})=1$ par définition donc il existe deux entiers $u$ et $v$. Par le théorème de Bézout on a:
 $$u\dfrac{a}{d}+v\dfrac{b}{d}=1$$
 autrement dit 
 $$ua+vb=d$$

\end{demo}

\begin{thm}

Soit $a$, $b_{1}$ et $b_{2}$ des entiers non nuls.

Si $a\wedge b_{1}=1$ et $a\wedge b_{2}=1$ alors $a\wedge (b_{1}b_{2})=1$  
\end{thm}

\begin{demo}

Si $a\wedge b_{1}=1$ il existe $u,v$ tels que $au+b_{1}v=1$

Si $a\wedge b_{2}=1$ il existe $u',v'$ tels que $au'+b_{2}v'=1$

donc: $(au+b_{1}v)(au'+b_{2}v')=1=a(auu'+b_{2}uv'+b_{1}u'v)+(b_{1}b_{2})(vv')$ donc $a\wedge (b_{1}b_{2})=1$  


\end{demo}

\begin{thm}[Théorème de Gauss]

Soit $a,b,c$ entiers avec $a, b$ non nuls.

Si $a$ divise $bc$ et $a$ premier avec $b$ alors $a$ divise $c$.

\end{thm}

\begin{demo}

Comme $a \wedge b=1$ il existe $u,v$ tels que $au+bv=1$ donc $cau+cbv=c$

$a$ divise $acu$ et $bcv$ donc $a$ divise $c$.

\end{demo}

\begin{thm}

Un entier $n$ divisible par $a$ et $b$ tels que $a\wedge b=1$ est divisible par $ab$.
\end{thm}

\begin{demo}

Si $a\mid n$ il existe un entier $k$ tel que $n=ak$.

Si $b\mid n$ alors $b\mid ak$ et comme $a\wedge b=1$, d'après Gauss, $b\mid k$. Il existe donc $k'$ tel que $k=bk'$ et on a alors $n=ak=abk'$ donc $n$ est divisible par $ab$

\end{demo}


\section{Nombres premiers}


\begin{defn}

Un entier naturel est dit \textbf{premier} s'il admet exactement deux diviseurs positifs (1 et lui-même)

\end{defn}

\begin{rem}
1 n'est pas premier.
\end{rem}

\begin{thm}

 Tout $n\in \mathbb{N}$ avec $n\geq 2$ admet au moins un diviseur premier.

 Si $n$ n'est pas premier et $n\geq 2$ alors il admet un diviseur premier compris entre 2 et $\sqrt{n}$
 \end{thm}


\begin{demo}

Si $n$ est premier, il admet bien un diviseur premier: lui-même.

Si $n$ n'est pas premier alors il admet un plus petit diviseur positif $p\neq 1$.
$p$ est premier sinon $p$ aurait lui-même un diviseur positif différent de 1 qui serait un diviseur de $n$, mais plus petit que $p$.

De plus, $n$ peut s'écrire $n=p \times r$ avec $p\leq r$ donc $p^{2}\leq p \times r$ soit $p^{2} \leq n$ et $p\leq\sqrt{n}$

\end{demo}

\begin{rem}
On utilise ce théorème de la manière suivante:

Si un naturel $n\geq 2$ n'admet pas de diviseur premier compris entre 2 et $\sqrt{n}$ alors $n$ est premier.

\end{rem}

\begin{exemple}

Pour savoir si 631 est premier, il suffit de tester tous les nombres premiers inférieurs à $\sqrt{631}\approx25,12$

\end{exemple}

\begin{thm}
Il existe une infinité de nombres premiers.
\end{thm}

\begin{demo}

On suppose qu'il existe un nombre fini de nombres premiers $p_{1}, p_{2}, .... ,p_{n}$


Soit $N=p_{1}\times p_{2}\times .... \times p_{n}+1$

$N$ n'est pas premier car pour tout $i$ de 1 à $n$, $N> p_{i}$.

 $N$ admet donc un diviseur premier dans la liste $p_{1}, p_{2}, .... ,p_{n}$. 

Soit $p_{k}$ ce diviseur premier de $N$.

$p_{k}$ divise $N$ et $p_{k}$ divise $p_{1}\times p_{2}\times .... \times p_{n}$ donc

$p_{k}$ divise $N-p_{1}\times p_{2}\times .... \times p_{n}$, soit $p_{k}$ divise 1: la seule possibilité est que $p_{k}=1$ qui est impossible car $p_{k}$ est premier.

\end{demo}

\begin{thm}

Tout entier naturel strictement supérieur à 1 admet une décomposition, unique à l'ordre des facteurs près, en produit de nombres premiers.

\end{thm}

\begin{exemple}

$72=2^{3}\times 3^{2}$
\end{exemple}


\begin{thm}[Petit théorème de Fermat]
Si $p$ est un nombre premier et si $a$ est un entier quelconque, alors $a^p-a$ est un multiple de $p$.
\end{thm}

\begin{lemme}
$p\vert C^k_p$ pour $k\in \{1,2,3,\cdots, p-1 \}$\\
\end{lemme}

\begin{demo}

$$C^k_p=\dfrac{p!}{(p-k)!k!}$$ 
$$C^k_p(p-k)!k!=p! $$
$$C^k_p(p-k)!k!=p\times (p-1)!$$

\end{demo}

\begin{lemme}
Soit $p$ un nombre premier $(k+1)^p \equiv 1+k^p \pmod p$
\end{lemme}

\begin{demo}
Par la formule du binôme on a,
$$(k+1)^p=\sum_{i=0}^p C^i_pk^i$$
de plus pour $i \in \{1,2,\cdots,p-1\}$ on a $p\vert C^i_p$ donc $ C^i_p\equiv 0 \pmod p$. 
Donc 
$$(k+1)^p \equiv 1+k^p$$
\end{demo}

\begin{demo}[Petit théorème de Fermat, par récurrence]
Initialisation : Pour $a=1$ on a $a^p=1^p=1$ donc l'assertion est vrai. \\

Hérédité : Supposons que maintenant que l'assertion est vrai pour $a=k$ alors montrons que l'assertion est vraie au rang $k+1$.\\

$(k+1)^p \equiv 1+k^p \pmod p$ or $k^p \equiv k \pmod p$ par hypothèse de récurrence. Donc $(k+1)^p \equiv 1+k \pmod p$.\\

Conclusion: pour tout entier $k$ et $p$ un nombre premier $k^p=k$.

\end{demo}


\end{document}

